\section{Reproduction}\label{sec:repro}

\paragraph{Version.} Ce rapport a été généré sur la branche \NGitBranch{} au commit \NGitHead. Étiquettes du dépôt :
\NGitTags.

\paragraph{Empreintes des prompts.} Analyse \NHashAnalyze, Plan \NHashPlan, \texttt{codes.py} \NHashCodes{} (combinée
\NHashCombined), Évaluer \NHashGrade, Reformuler \NHashReformulate, Vérifier \NHashVerify.

\paragraph{Environnement.} Un seul environnement conda, dont les versions exactes sont figées dans
\texttt{requirements.lock.txt}. Les figures sont dessinées par un interpréteur disposant de matplotlib (variable
\texttt{FIG\_PYTHON}), à partir de \texttt{report/data.json} uniquement.

\paragraph{Commandes.}
\begin{lstlisting}
python evaluate.py --split dev --temperature 0
python agent/eval_plan_vs_raw.py --split dev
python agent/eval_loop.py --split dev
python agent/eval_loop.py --split dev --start g
python eval/pdf_truth.py
python eval/audit_chunks.py
python eval/show_article.py PenalFr 218
python eval/verifier_calibration.py --n 24
python eval/noise_check.py
python eval/pdf_factcheck.py results/noise_check.json
python report/build_report.py
\end{lstlisting}

\section{Incidents et leçons}\label{sec:incidents}

\paragraph{Écart d'environnement.} Une exécution du script de comparaison a donné, pour la variante A1 sur le
développement, 77,1~\% de Hit@6 au lieu des \NVoneoneDevHitsix~\% attendus \averifier. Avec le code engagé et
l'environnement figé, la valeur attendue a été reproduite à l'identique, question par question. La cause de l'écart n'a
pas pu être établie à partir des fichiers ; l'hypothèse retenue est une exécution dans un autre environnement \averifier.
Leçon : figer l'environnement et enregistrer la provenance de chaque exécution.

\paragraph{HEAD détaché et paramètre perdu.} Un passage sur un commit ancien en HEAD détaché a remplacé le fichier de
récupération (perte du paramètre de filtre doux) et les prompts gelés par leurs versions antérieures, puis une étiquette
a été posée sur ce commit. L'état a été récupéré en revenant sur la branche de travail et en y reportant la seule
modification utile ; le commit détaché est conservé dans la branche \NGitRescue{} et l'étiquette a été déplacée.
Leçon : vérifier la branche avant chaque commit, et dater les résultats par l'empreinte de leurs prompts.

\paragraph{Jeu de test vu deux fois.} Le jeu de test a servi à la mesure finale de v1.1, puis à la comparaison des
variantes de planification (\texttt{agent/LEARNING.md}). Il ne peut plus confirmer les résultats de l'agent.

\paragraph{Hypothèse sur l'arabe.} Une hypothèse sur l'avantage de certaines langues a été formulée puis retirée faute de
mesure suffisante \averifier.

\section{Structure du dépôt}
Fichiers et dossiers suivis à la racine :
\begin{lstlisting}
.gitignore
L1Context.drawio
L2Flow.drawio
L3_data.md
PROJECT_TECHNICAL_EXPLANATION.txt
README.md
TODO.txt
agent
app.py
config.py
data
docs
eval
eval_harness.py
evaluate.py
loaders.py
out.json
pipeline.py
preprocessing.py
requirements.lock.txt
reranker.py
results
retriever.py
splitter.py
translate.py
\end{lstlisting}

\begin{itemize}
  \item \texttt{agent/} : nœuds de l'agent, graphe, scripts d'expériences ;
  \item \texttt{eval/} : questions, découpage, vérité de terrain PDF, audit, étalonnage, harnais d'évaluation ;
  \item \texttt{results/} : sorties des exécutions (source de tous les nombres) ;
  \item \texttt{report/} : ce rapport et ses scripts de génération.
\end{itemize}
